\subsubsection{Accessing S3 Buckets}

\newline
OpenScienceLab user servers typically run in AWS region \texttt{us-west-2}, which houses all public NASA data. This provides cheap, low-latency data access.

You may access any public AWS S3 bucket or private bucket to which you have permission with either the \texttt{aws-cli} (command line interface) or \texttt{boto3} (Python package).

\begin{framed}
\textbf{Configure the AWS Client in your OpenScienceLab account for non-public bucket access.}\\
\textit{Prerequisite: an AWS access key that provides access to the S3 bucket}

\begin{enumerate}
\item Open a terminal
\item Run the command: \texttt{aws configure}
\item When prompted, provide:\begin{itemize}
\item \texttt{AWS Access Key ID}
\item \texttt{AWS Secret Access Key}
\item \texttt{Default region name}: ``us-west-2''
\item \texttt{Default output format}: ``json''
\end{itemize}
\end{enumerate}
\end{framed}

\begin{framed}
\textbf{Warning}\\
If accessing a public S3 bucket anonymously, you must either:

\begin{itemize}
\item Provide the \texttt{--no-sign-request} argument to AWS-cli
or
\item Provide an \texttt{UNSIGNED} config to the boto3 S3 bucket resource:\begin{itemize}
\item \texttt{s3 = boto3.resource('s3', config=Config(signature\_version=UNSIGNED))}
\end{itemize}
\end{itemize}
\end{framed}

\paragraph{Example Notebook Demonstrating S3 Bucket Access with \texttt{AWS-cli} and \texttt{boto3}}

\href{https://github.com/ASFOpenSARlab/NISAR-Early-Adopters-Workshop-Oct-2024/blob/main/S3\_Buckets\_Upload\_Download.ipynb}{S3\_Buckets\_Upload\_Download.ipynb}

\begin{verbatim}

\end{verbatim}